\subsection{The finite Levi comparison for the actual zero-Gram operator}
\label{sec:flag-finite-levi}

We complete the two finite interfaces in (GA9), retaining the
original \(\mathrm{Full}\ X\), column representatives, and
\(\mathrm{TypedRed}\) throughout. The group in this section is
\[
 \mathcal G=\prod_{v\in\mathrm{TR.Pl}}\mathrm{GL}_k(\mathbb F_{q_v}),
 \qquad U=\prod_v U_k(\mathbb F_{q_v}),
\]
where \(U_k\) denotes the upper unitriangular group. Let \(p(x)\)
be the product of the original adjacent-entry characters on \(U\),
extended by zero to \(\mathcal G\). Each additive character is
nontrivial, by the \(\mathrm{psi\_ne}\) field of \(\mathrm{Full}\).
Define
\[
 \gamma(x)=|U|^{-1}\sum_{y\in\mathcal G}p(yxy^{-1}),
 \qquad \Gamma_1=\gamma(1)=|\mathcal G|/|U|.
 \tag{FL1}
\]
Thus \(\gamma=\mathrm{gamProd}\) and \(\Gamma_1=\mathrm{Gam1}\);
the finite group is not a replacement for the column space.

\begin{lemma}[Finite Levi Gram bound]
\label{lem:flag-finite-levi-gram}
The original matrix \(GG(i,j)=\overline{\gamma(j^{-1}i)}\) satisfies
\[
 0\le GG\le P_g\Gamma_1\,1,\qquad
 P_g=\prod_v(1+3/q_v)\le2
 \tag{FL2}
\]
under the original standing size conditions.
\end{lemma}
\begin{proof}
First use a single field \(\mathbb F_q\). For a unitary
representation \(\pi\) of \(G=\mathrm{GL}_k(\mathbb F_q)\),
\(|U_k|^{-1}\sum_{u\in U_k}p(u)\pi(u)\) is the orthogonal
projection onto the inverse-character subspace. Its sum of
conjugates over \(G\) is positive, central, and equals
\(\sum_x\gamma(x)\pi(x)\). On an irreducible of dimension \(d\)
it acts by \(|G|m/d\), where \(m\) is the multiplicity of
that character.

For completeness, the multiplicity-one argument here is the
elementary finite Gelfand--Graev argument. Bruhat decomposition
gives representatives \(tw\) for the \(U_k\)-double cosets, with
\(t\) diagonal and \(w\) a permutation matrix. A double coset
supporting a nonzero bi-character function must satisfy the
character-matching condition on \(U_k\cap twU_k(tw)^{-1}\).
If two adjacent positions \(j,j+1\) in \(w\) have increasing
values, conjugate the simple root \(1+aE_{j,j+1}\) by \(tw\).
Its image is an upper root. Matching for all \(a\) forces
the image root to be simple, and the ratio of its two diagonal
entries of \(t\) to be one. Indeed a nonsimple root has trivial
character, and a nontrivial additive character cannot satisfy
\(\psi(ca)=\psi(a)\) for all \(a\) unless \(c=1\).
Every increasing run of the permutation therefore consists of
consecutive integers, and \(t\) is constant on its value interval.
The intervals are disjoint, their successive transitions are
descents, and hence they occur in reverse interval order.
The monomial matrix \(tw\) consequently consists of scalar identity
blocks in reversed block order. It is fixed by the anti-involution
\(x\mapsto w_0 x^{\mathsf t} w_0\), where \(w_0\) reverses all
coordinates. This anti-involution preserves \(U_k\) and the
adjacent-entry character. It fixes every bi-character function:
on each supported double coset this follows from its fixed
representative and the two identical character factors, and on
the other cosets from the vanishing condition. Convolution is
reversed by the anti-involution, so the character Hecke algebra
is commutative. Since over \(\mathbb C\) it is the endomorphism
algebra of the induced character representation (up to taking the
opposite algebra), semisimplicity and Frobenius reciprocity give
\(m\le1\). The argument also applies to the inverse character.

A representation admitting that character has dimension at least
\(\prod_{a=1}^{k-1}(q^a-1)\). Here is the usual mirabolic induction
without a character-degree formula. Restrict to the group
\(\mathrm{GL}_{k-1}\ltimes\mathbb F_q^{k-1}\) fixing the last
row. A character vector lies in the nonzero weight space of the
last-column vector group defined by the last coordinate. All its
\(q^{k-1}-1\) nonzero characters have equally large weight spaces,
by transitivity of \(\mathrm{GL}_{k-1}\). The stabilizer of the
chosen character in \(\mathrm{GL}_{k-1}\) is the smaller
mirabolic group; its chosen weight space contains the same
adjacent-root character for \(U_{k-1}\). Repeat within that
weight space, starting the induction at dimension one. This
argument does not require irreducibility of the restrictions.
Using \(|U_k|=q^{k(k-1)/2}\), for \(m=1\) we have
\[
 \frac{|U_k|}{d}\le
 \prod_{a=1}^{k-1}(1-q^{-a})^{-1}
 \le \left(1-\frac1{q-1}\right)^{-1}
 \le 1+\frac3q\quad(q\ge3).
\]
The middle inequality follows by multiplying numbers \(1-x_a\)
with \(0\le x_a\le1\) and bounding their product below by
\(1-\sum_a x_a\). For \(m=0\) the convolution is zero.
Applying this spectral argument to the regular representation,
then taking tensor products over the actual places, proves the
matrix inequality in (FL2); transpose or complex conjugation of a
Hermitian positive matrix leaves its real spectral bounds intact.

To check the number of factors, use the original coefficient
reduction. At the two completely split rational level primes,
\(\mathcal O_E/g\mathcal O_E\) consists of \(f\) pairs of copies
of each of \(\mathbb F_{g_0}\) and \(\mathbb F_{g_1}\),
exchanged by conjugation. The original idempotent \(e_w\) and
its conjugate are orthogonal and sum to one modulo \(g\)
(the faithful reduced lattice detects this equality).
Thus \(e_w\) selects one factor per pair. The joint surjectivity
and kernel equality of \(\mathrm{TypedRed}\) identify its
product of residue fields with exactly this selected ring, as
in the row-congruence proof. An isomorphism of these products
permutes primitive field idempotents and preserves their
characteristics. Hence each \(g_i\) occurs exactly \(f\) times,
not an unspecified number of times. The size conditions
\(g_i\ge(kb)^8,\ f\le k,\ b\ge2,\ k\ge32\) now give
\(\log P_g\le3f/g_0+3f/g_1\le6k/(2k)^8<\log2\).
\end{proof}

\begin{lemma}[The original columns and their row masses]
\label{lem:flag-finite-columns}
For each original row \(\Lambda\), the original finite
column representatives in that row give a disjoint decomposition
\[
 \mathcal G=\coprod_{\rho:\,\lambda(\rho)=\Lambda}
                   A_\Lambda s_\rho U
 \tag{FL3}
\]
with unique triples \((a,\rho,u)\).
Here \(s_\rho=\mathrm{scOf}(\rho)\) and
\(A_\Lambda=\mathrm{Aset}(\Lambda)\) are the actual reductions.
In particular,
\[
 |A_\Lambda|^{-1}=w_{\rm row}(\Lambda),\qquad
 d_u=d_\pi\sum_\Lambda\#\{\rho:\lambda(\rho)=\Lambda\}
     =d_\pi\Gamma_1 M .
 \tag{FL4}
\]
\end{lemma}
\begin{proof}
Let \(H_u\) be the pointwise stabilizer of the original rational
vectors \(e_a,f_a,z_u\). Its finite and rational points are
precisely the groups \(\mathrm{Hfin}\) and \(\mathrm{Gam}\)
in the column dictionary. A representative \(h_\rho\) defines
the row lattice by
\((h_\rho x_u)^{-1}v\in\widehat L_0\). The given \(x_u\) preserves
\(\widehat L_0\), so this means
\(v\in h_\rho\widehat L_0\). Local-global equality for integral
lattices now identifies its finite completion with that lattice.
Since \(h_\rho\) fixes both hyperbolic frames, the splitting
of \(L_0\) is also a splitting of every such row, with the same
isotropic and dual ideals. Only the middle lattice varies.
In particular its middle quotient \(X^\perp/X\) identifies
integrally with the middle lattice. The rational \(H_u\)-stabilizer
of the row therefore identifies, in both directions, with the
isometries of this quotient fixing the class of \(z_u\): extend
such an isometry by the identity on the two hyperbolic frames.
This is the automorphism group in the original row weight.
It is finite because it is integral and lies in a compact
product at the real places on the positive complement of \(z_u\).

For the chosen original base column \(h_\Lambda\) of a row,
all its finite representatives are \(h_\Lambda K_H\), where
\(K_H=H_u(\mathbb A_f)\cap K\) and \(K\) is the stabilizer of
\(\widehat L_0\). Their relative matrices, in the level
coordinates, lie in \(x_u^{-1}K_Hx_u\). We show that reduction
on this group is onto \(\mathcal G\), and that its inverse
image of \(U\) is the level subgroup. This is the needed
column surjectivity, not strong approximation of a rational
unitary group.

At a place dividing \(g\) the actual different is a unit,
and \(T=-\langle z_u,z_u\rangle\) is a unit, since the original
word norm is prime to both level primes. In the \(x_u\)
coordinates the fixed \(e_a,f_a,x_u^{-1}z_u\) give an integral
unimodular hyperbolic summand and a unimodular negative line.
Their orthogonal complement is self-dual of rank \(k\).
The CM extension is split at the place. Thus the two components
of this complement are free modules of rank \(k\) over
\(\mathcal O_{F,v}\), in perfect duality. Its stabilizer is
\(\mathrm{GL}_k(\mathcal O_{F,v})\), acting by inverse transpose
on the other component. The original \(e_0,f_0,z_0\) agree
with the fixed vectors modulo \(g\) by their pinned
approximations (denominators prime to \(g\) are units here).
The original \(bw,bwd\) give dual bases of the selected
and opposite residual complement: this follows from the
stated Gram congruences and the generation of the reduced
lattice by these vectors and the fixed frames. Projection
onto the exact complement changes these vectors by zero
modulo \(g\). In these bases the coefficients
\(\langle bwd_i,\kappa bw_j\rangle\) are exactly the selected
matrix entries of reduction. The determinant is a unit, so
the fallback branch of \(\mathrm{levi}\) is not used. Any
invertible residue matrix lifts to an invertible integral
matrix, with the inverse-transpose action on the dual and
identity on the fixed summands. Choose these lifts at the
finitely many places dividing \(g\) and the identity elsewhere;
this proves simultaneous surjectivity for the actual relative
column group.

The defining \(K_2\) congruences for the fixed \(e_0,z_0\) are
automatic on that group. The remaining congruences require
\(bw_j\) to be fixed modulo its preceding \(bw_i\)'s and the
isotropic span. They are exactly upper unitriangular reduction.
The dual action is forced by unitarity; outside \(g\) there
is no further congruence. The inverse image is therefore
the original \(x_u^{-1}K_ux_u\). In particular the relative
matrix from the chosen base to column \(\rho\) reduces
to \(s_\rho\), as defined in the source.

The rational stabilizer of the row maps to \(\mathcal G\) by
the same conjugation and reduction; its image is \(A_\Lambda\).
This map is injective, and the image acts freely on
\(\mathcal G/U\). Indeed if a finite-order stabilizer element
has reduction conjugate into \(U\), its action on each
whole reduced lattice is unipotent: it fixes the hyperbolic
frames and last vector and acts dually on the other component.
At a residue characteristic \(p\), a finite-order integral
matrix with unipotent reduction has \(p\)-power order. To see
this, raise it to a \(p\)-power to kill the residual unipotent
part and leave only its prime-to-\(p\) order. For a matrix
\(B\equiv1\) of prime-to-\(p\) order \(n\),
\(1+B+\cdots+B^{n-1}\) reduces to the invertible scalar \(n\).
Thus \(B^n-1=0\) implies \(B=1\). Applying this at both
distinct rational characteristics forces the original element
to be the identity. This proves both assertions and gives
\(|A_\Lambda|^{-1}=w_{\rm row}(\Lambda)\).

The actual \(\mathrm{ColReps}\) assertion gives exactly one
representative per rational-left and \(K_u\)-right double
coset. For two representatives defining the fixed row,
a rational left factor relating them must preserve that row
because the right factor preserves the completed lattice.
Conversely such factors give precisely the double cosets
in the row. By the surjective reduction and its just proved
inverse-image property these are the
\(A_\Lambda\backslash\mathcal G/U\) cosets. Thus (FL3) is
a bijection with the original representatives, not a
replacement column count. Their weight is the displayed
reciprocal order, so their number is
\(\Gamma_1w_{\rm row}(\Lambda)\).
Finally the original column representation is specified by
arbitrary vectors of the compact type at the column
representatives: the rational stabilizer at a column is
trivial by the same two-characteristic argument. Evaluation
therefore gives the full column dimension, also in the
given Riesz identification. This proves (FL4).
\end{proof}

\begin{lemma}[The original zero-Gram norm expansion]
\label{lem:flag-finite-norm-expansion}
In the absolutely convergent original diagonal expansion, let
\(K''\) denote its nonidentity zero-Gram part. With the
amplitudes and weights of (GA3), one has
\[
 \|K''\|_{\mathrm{HS}}^2
 \le P_g\Gamma_1\!
  \sum_{(t,t')\in\mathrm{allPairs}(\mathrm{degSet})}a_ta_{t'} .
 \tag{FL5}
\]
\end{lemma}
\begin{proof}
For a pair of rows \(Q=(\Lambda,\Lambda')\), let
\(Z_Q(i)\) be the sum of the original \(\mathrm{ampN}(t)\)
operators over zero-Gram nonidentity terms with row
\((\Lambda,\Lambda',i)\), viewed in their Hilbert--Schmidt
space. These sums exist by the original absolute convergence.
In this notation \(\mathrm{vec}(t)\) is the corresponding
single-row vector of weight
\(\sqrt{w_{\rm row}(\Lambda)w_{\rm row}(\Lambda')}\).

We first check the finite coefficients of the actual column
blocks. The normalized base return \(\kappa_t\) preserves the
integral lattice: the return takes the source row lattice to
the target row lattice. The original admissibility at \(g\)
makes its reduction on the isotropic flag upper unitriangular
and on the last-vector class the identity (the integral
congruence proof of (FC7)). It therefore acts on the selected
orthogonal quotient of rank \(k\) by the actual Levi
matrix \(i(t)\). The adjointness of the complementary
residue components forces its action on the other component.
Multiplying by the two relative column matrices (inverse on
the target) replaces \(i(t)\) by
\(s_\rho^{-1}i(t)s_{\rho'}\); this is ordinary composition
of the induced maps on the quotient. The fixed-frame
congruences are unchanged. In particular the resulting
return belongs to \(K_2\) precisely when this matrix
belongs to \(U\).

The original \(\chi_N\) is the finite character of the
sum of \(\langle f_{0,a},\kappa\,\mathrm{nxt}(e_0,z_0)_a\rangle\),
divided by \(g\). Both relative column matrices fix the
\(e_0,f_0,z_0\) modulo \(g\); by unitarity they leave
each of these pairings unchanged modulo \(g\). All matrices
and pairings in this computation are integral, and outside
the level \(g\) is a unit. Hence the finite character is
unchanged, exactly, not up to an additional scalar.
The remaining factor \(\chi_L\) uses the adjacent entries
\(\langle bwd_j,\kappa bw_{j+1}\rangle\), divided by \(g\).
These coefficients lie in the selected \(e_w\) component,
and the given \(\mathrm{psi\_eq}\) dictionary identifies the
character with the product of the prescribed residue
characters. Thus with the zero outside \(K_2\) the factor
is precisely \(p(s_\rho^{-1}i(t)s_{\rho'})\).
This proves the original \(\mathrm{ColPhase}\) and
\(\mathrm{ChiNIndep}\) identities needed for these blocks.
It follows directly from the original level-character
splitting and \(\mathrm{ampN}\) definition that the
\((\rho,\rho')\)-block of \(K''\) is
\[
 B_{\rho,\rho'}=\sum_{i\in\mathcal G}
               p(s_\rho^{-1}i s_{\rho'})\,Z_Q(i).
 \tag{FL6}
\]

Next consider \(a\in A_\Lambda,\ b\in A_{\Lambda'}\).
They have rational lifts \(\alpha,\beta\) stabilizing their
row lattices, fixing the original two hyperbolic frames and
last vector. Multiplication of a rational return by these
lifts, \(\gamma\mapsto\alpha\gamma\beta^{-1}\), is a
bijection on the corresponding original double cosets:
the lifts normalize the respective arithmetic unipotent
groups. They preserve the admitted-frame condition, the
Gram-zero condition and identity membership. In detail,
the paired frames are transported by an isometry fixing
one flag, and the other lift fixes the second reference
flag pointwise; integral membership and admission are
unchanged since the row lattices are preserved. Relevance
is likewise invariant since the differential root characters
are preserved by conjugation with an isometry fixing the
two hyperbolic frames and last vector. This also follows
directly by transporting the relevant common stabilizer
and its two characters. The Levi label is changed to
\(aib^{-1}\).

The \(\chi_N\) character is unchanged by these relative
matrices, by the preceding fixed-frame calculation.
In the real integral defining \(\mathrm{OO}_t\), transport
the stabilizer fundamental set by \(\alpha\) and change
variables by that same element. The two original functions
are equivariant under \(H_u(\mathbb R)\) with its given
unitary compact-type action. Consequently the resulting
operator differs from the original by left and right
unitary actions of \(\alpha,\beta\). Both the fundamental
set transport and replacement by the chosen double-coset
representative are valid by the already established
absolute convergence and phase compatibility of the
orbital expansion. The induced unitary action on the
Hilbert--Schmidt space is independent of the term.
Summing on a label fibre gives, in particular, the
\(\mathrm{SatGram}\) identity
\[
 \langle Z_Q(aib^{-1}),Z_Q(ajb^{-1})\rangle
       =\langle Z_Q(i),Z_Q(j)\rangle.
 \tag{FL7}
\]
This uses saturation of the actual zero-Gram nonidentity
set, not an arbitrary selection of its terms.

Now expand the norm over the genuine column blocks.
The character \(p\) is left and right \(U\)-equivariant,
including on its zero set. Using (FL7), the squared norm
of \(\sum_i p(x^{-1}iy)Z_Q(i)\) is unchanged on replacing
\(x,y\) by \(a x u,b y u'\), respectively. The factors
from \(u,u'\) have modulus one, and (FL7) gives equality
after the simultaneous change of label. Thus (FL3)
and the Hilbert--Schmidt block identity give the
contribution of the pair \(Q\) to the squared norm as
\[
 \frac{1}{|A_\Lambda||A_{\Lambda'}||U|^2}
  \sum_{x,y\in\mathcal G}
      \left\|\sum_i p(x^{-1}iy)Z_Q(i)\right\|_{\rm HS}^2 .
\]
For fixed \(i,j\), the coefficient of the inner product
(conjugate-linear in the first argument) in this sum before
division is
\[
 \begin{aligned}
 \sum_{x,y}\overline{p(x^{-1}iy)}p(x^{-1}jy)
 &=|U|\sum_y p(y^{-1}i^{-1}jy)\\
 &=|U|^2\gamma(i^{-1}j).
 \end{aligned}
\]
In the first equality substitute \(z=x^{-1}iy\); only
\(z\in U\) contributes and the character factors cancel.
Since \(\gamma\) is central and
\(\gamma(s^{-1})=\overline{\gamma(s)}\),
this is \(|U|^2 GG(i,j)\), with precisely the original
orientation. Consequently, if \(V_i\) is the direct sum
of the vectors
\(\sqrt{w_{\rm row}(\Lambda)w_{\rm row}(\Lambda')}Z_Q(i)\),
we obtain the actual norm expansion
\[
 \|K''\|_{\rm HS}^2=\sum_{i,j}GG(i,j)
                             \langle V_i,V_j\rangle .
\]
The spectral bound (FL2) applied in each Hilbert-space
coordinate is at most
\(P_g\Gamma_1\sum_i\|V_i\|^2\).
Expanding each row-label sum by the triangle inequality
gives (FL5), with both row weights, not merely a sum
over the Levi labels.
\end{proof}

\begin{theorem}[The diagonal zero-Gram saving on the matched local data]
\label{thm:flag-zero-gram-saving}
Under the standing hypotheses and the matched local-data hypothesis
of the global amplitude theorem, the original diagonal
nonidentity zero-Gram operator satisfies
\[
 \boxed{\displaystyle
    \|K''_{uu}\|_{\rm HS}^2\le d_u\exp(-kf\log k).}
 \tag{FL8}
\]
In particular the full zero-Gram summation and finite norm
comparison have no further arithmetic or counting hypothesis.
The existence of the required matched local data from the
original standing hypotheses is a separate remaining interface.
\end{theorem}
\begin{proof}
Use (GA7), (FL4) and (FL5):
\[
 \|K''_{uu}\|_{\rm HS}^2
 \le 2\Gamma_1d_\pi M\,e^{-3kf\log k/2}
 =2d_u e^{-3kf\log k/2}
 \le d_u e^{-kf\log k}.
\]
The final inequality uses \(k\ge32,\ f\ge1\).
All sums refer to the original full zero-Gram nonidentity set.
Its amplitude-pair summability is part of (GA7), and the
operator sum is the given absolutely convergent expansion.
No assertion concerning generic nonidentity returns or the
final arithmetic theta comparison follows from this theorem.
\end{proof}
